\section{部署与监控实践}
\subsection{CI/CD 集成}
Graham 建议把 Agent 与 SWE-bench 结合到 CI 流水线：每天或每个 commit 运行 Lite 集合，并把失败样本生成 issue 供模型团队研究。自动化脚本可以在多个环境中并行执行 Lite/Full，并上传日志供人类审查。

\begin{lstlisting}[language=bash]
#!/usr/bin/env bash
set -ex
cd /workspace/swe-bench
python run_swe_bench.py --split lite --model codegen
python run_swe_bench.py --split full --tests smoke
tar czf swe-bench-results.tar.gz outputs/
gzip -c outputs/logs/latest.log > swe-bench-results/logs/latest.log.gz
notify-team --subject "SWE-bench run finished" --files swe-bench-results.tar.gz
\end{lstlisting}

脚本中的 \texttt{notify-team} 可以接入 Slack/Email，把失败 diff 以及通用输入输出一起发送。

\begin{knowledgebox}{CI/CD 运行策略}
把 Agent 运行拆成“预热、执行、回放”：预热阶段重建 repo 与依赖，执行阶段运行 Agent，回放阶段把日志 + diff 上传供开发者复审。
\end{knowledgebox}

\subsection{小规模实验}
在将 Agent 扩大部署前，先在 \texttt{sandbox} 仓库做 1-2 个试点任务，用 Agent 生成 patch，但只在内网 review，并保持每轮运行在 15 分钟以内，确保可控。

\begin{importantbox}{短期实验建议}
保留人工检查点：只允许 Agent 修改少于 3 个文件，所有操作必须在 \texttt{git} 里打 patch 后由人类确认。
\end{importantbox}

\begin{warningbox}{避免主分支污染}
不要让 Agent 直接 push 到 `main`/`master`；应该在 feature 分支生成 patch，待 review 完成后再合并。
\end{warningbox}

\subsection{本章小结}
\begin{itemize}
    \item 把 Agent 嵌入 CI 后能定期捕捉 regressions，并把失败 trace 回 feed。
    \item 小规模的内网实验可以降低迭代风险并积累经验。
\end{itemize}

